\section{The finite verification of the minimum at level two}\label{sec:minimum}
The minimum at strict level one is $m_1=8$~\cite{DEFH2023}.
This section explains the verification that $m_2=12$. The calculations
exclude the finite set of eleven-vertex extensions left by the
mathematical reduction below. The twelve-vertex graph $H_2$ from
Section~\ref{sec:family4} supplies the matching example.

\subsection{Why these extensions suffice}
First let $k\geq2$ and choose a graph $J$ of strict level $k$ and order $m_k$.
If $v$ is the first vertex of a level-$k$ reduction, then $J-v\in\D_{k-1}$
by minimality. The graph $X=L_J(v)$ lies in $\D_{k-1}\setminus\D_{k-2}$:
otherwise that first deletion would put $J$ in $\D_{k-1}$. Since $J$ is not
a cone, $v$ has a non-neighbour. Counting $X$, $v$ and that non-neighbour gives
\[
 m_k\geq m_{k-1}+2\qquad(k\geq2).
\]
If a graph of strict level two had at most eleven vertices, we could choose
$J$ above with $k=2$ and $|V(J)|\leq11$. Put
$U=V(J)\setminus N_J[v]$. Since $m_1=8$ and $U$ is non-empty,
the only possible sizes are
\[
 (|V(X)|,|U|)=(8,1),\ (9,1),\ (8,2).
\]
We use the following small-order classification~\cite{DEFH2023}:
every non-empty graph of order at most eight whose clique complex is acyclic
over $\F$ belongs to $\D_1$; exactly seven isomorphism classes among these
graphs lie outside $\D_0$, all of order eight.

If $U=\{u\}$, put $Z=N_J(u)\subseteq V(X)$. Fold each vertex of
$V(X)\setminus Z$ onto $v$. The remainder is the complete join
$X[Z]\star\overline{K_2}$, where the two new vertices are $u,v$.
By Lemmas~\ref{hom:lem:acyclic} and~\ref{hom:lem:fold}, its clique complex
is acyclic. The suspension identity~\eqref{hom:eq:suspension} then gives
acyclicity of $\Cl(X[Z])$, including degree $-1$; in particular $Z$ is non-empty.
If $Z=V(X)$, we already have $X[Z]\in\D_1$. Otherwise $|Z|\leq8$, so the
small-order classification gives the same conclusion. Lemma~\ref{ker:lem:join}
now puts the remainder in $\D_1$, and prepending the folds puts $J$ in $\D_1$,
a contradiction. Both cases with $|U|=1$ are therefore excluded without
enumerating nine-vertex graphs.

Only $(|V(X)|,|U|)=(8,2)$ remains. Since $X\in\D_1\setminus\D_0$, it is
one of the seven graphs in the classification. This is why the following
seven seeds, rather than a catalogue of all eleven-vertex graphs, suffice.

\subsection{Seven seeds with explicit reductions}
Use labels $0,\ldots,7$ and specify each seed by its missing edges in
Table~\ref{min:tab:seeds}. The notation $ab$ means the unordered pair $\{a,b\}$.
\begin{table}[htbp]
\centering\small
\caption{The seven eight-vertex seeds, specified by their missing edges.}\label{min:tab:seeds}
\begin{tabular}{cl}
\toprule
Graph & Missing edges\\\midrule
\(X_1\)&\(01,02,07,13,16,25,26,35,45,46,47\)\\
\(X_2\)&\(02,03,05,13,14,24,26,34,57,67\)\\
\(X_3\)&\(01,04,05,17,24,26,35,36,37,56\)\\
\(X_4\)&\(01,04,05,17,24,26,27,35,36,56\)\\
\(X_5\)&\(02,03,04,12,15,16,35,37,46,67\)\\
\(X_6\)&\(01,02,12,15,26,34,37,46,57\)\\
\(X_7\)&\(02,03,05,12,14,16,34,57,67\)\\
\bottomrule
\end{tabular}

\end{table}
None has a dominated vertex. In Table~\ref{min:tab:reductions}, the first
vertex has a four-vertex path as its neighbourhood, so it may be deleted at
level one. The subsequent folds leave the indicated cone. Each seed
therefore has strict level one.
\begin{table}[htbp]
\centering\small
\caption{First deletions and folds of the seven seeds. Paths list the
neighbours in path order; folds take place after the first deletion.}\label{min:tab:reductions}
\begin{tabular}{ccccc}
\toprule
Graph & First vertex & Path & Subsequent folds & Apex\\\midrule
\(X_1\)&\(0\)&\(4,3,6,5\)&\(4\to2\)&\(7\)\\
\(X_2\)&\(3\)&\(7,2,5,6\)&\(2\to1,\ 7\to0\)&\(6\)\\
\(X_3\)&\(0\)&\(3,2,7,6\)&\(3\to1,\ 2\to5\)&\(4\)\\
\(X_4\)&\(0\)&\(2,3,7,6\)&\(2\to1\)&\(4\)\\
\(X_5\)&\(0\)&\(1,7,5,6\)&\(1\to4\)&\(2\)\\
\(X_6\)&\(1\)&\(3,6,7,4\)&\(3\to5,\ 2\to4\)&\(0\)\\
\(X_7\)&\(0\)&\(1,7,4,6\)&\(1\to3\)&\(2\)\\
\bottomrule
\end{tabular}

\end{table}
Their edge counts are $17,18,18,18,18,19,19$. For $X_2,\ldots,X_5$, the
vertices of degree four induce graphs with degree sequences
$(1,1,1,1)$, $(0,1,1,2)$, $(1,1,2,2)$ and $(2,2,2,2)$, respectively.
The two degree-four vertices are non-adjacent in $X_6$ and adjacent in $X_7$.
These invariants distinguish all seven classes.
\runcheck{small-orders}
The calculation verifies these reductions and invariants. It also reads all
$13\,599$ graphs in the supplied catalogues of orders zero through eight,
computes augmented reduced homology over $\F$, and tests membership in
$\D_0$ and $\D_1$. It finds exactly the seven classes above and records an
explicit isomorphism for each. The empty graph is included and has
$\widetilde\beta_{-1}=1$; it is not counted as acyclic.
The catalogue sources and file identities are recorded in
\path{THIRD_PARTY.md} and \path{data/catalogues/provenance.json}.
The main command \texttt{python reproduce.py} also checks completeness
independently: it tests pairwise non-isomorphism and compares the number of
entries at each order with Burnside's count of the orbits of edge sets under
all vertex permutations.

\subsection{All choices of the two external neighbourhoods}
For each $X_i$, add $v,u_1,u_2$. The neighbours of $v$ are precisely
$V(X_i)$; the vertices $u_1,u_2$ may have arbitrary neighbourhoods
$Z_1,Z_2\subseteq V(X_i)$, and their mutual edge may be present or absent.
The program considers all $2\cdot256^2=131\,072$ ordered choices for each
seed. It also records those with $Z_1\leq Z_2$, ordering subsets by their
binary encodings. Exchanging $u_1,u_2$ identifies the two orders, leaving
$2\binom{257}{2}=65\,792$ choices per seed and $460\,544$ in total.
These are choices of extension data; different choices can produce
isomorphic graphs.

For each extension, a non-zero $q$ excludes every dismantling class by
Lemma~\ref{lem:qobstruction}. If $q=0$, the program applies the exact
recursion~\eqref{eq:first} at levels zero, one and two. Every possible first
deletion is considered when necessary, and each recursive call has fewer
vertices. A zero $q$ alone never establishes membership.
\begin{table}[htbp]
\centering\small
\caption{The $65\,792$ extension choices per seed, after exchanging the two
external vertices. The four columns partition the choices; the last column
has $q=0$ and lies outside $\D_2$.}\label{min:tab:extensions}
\begin{tabular}{crrrr}
\toprule
Graph & \(q\ne0\) & \(\D_0\) & \(\D_1\setminus\D_0\) & \(q=0,\ \notin\D_2\)\\\midrule
\(X_1\)&32823&27320&3610&2039\\
\(X_2\)&36048&25593&1821&2330\\
\(X_3\)&36283&25130&1920&2459\\
\(X_4\)&35171&26246&2219&2156\\
\(X_5\)&32323&28668&2930&1871\\
\(X_6\)&36573&25833&1331&2055\\
\(X_7\)&33543&28579&1733&1937\\
\bottomrule
\end{tabular}

\end{table}
\runcheck{order-eleven}
The program evaluates all $917\,504$ ordered choices and compares the
classification of all $460\,544$ choices up to exchange with a separate
implementation using sets of vertices and clique enumeration. The report gives the
ordered counts, the counts in Table~\ref{min:tab:extensions}, and the counts
with $Z_1=Z_2$. For each category it checks
$N_{\rm ordered}=2N_{\rm unordered}-N_{\rm equal}$.
No extension has strict level two. Combined with the preceding case
reduction, this proves that no graph of order at most eleven has strict level two.

The seven-root certificate in Section~\ref{sec:family4} proves that $H_2$
has strict level two and twelve vertices. The command below is an additional
check of the article's direct verification, not a further requirement for
this proof. It uses the automorphism
$(4\,5)(6\,7)(8\,9)(10\,11)(12\,13)$ of $H_2$ to group its neighbourhoods.
Four representative neighbourhoods fold to the octahedral graph
$K_{2,2,2}$, whose non-zero $q$ excludes ordinary dismantling; the other
representatives have no dominated vertex. The calculation checks these
obstructions and the positive reduction.
\runcheck{h2}
It verifies a level-two reduction and excludes a first level-one deletion
at every vertex, retaining the original labels $2,\ldots,13$.
Thus the upper and lower bounds give $m_2=12$.
The descent inequality proved above therefore gives $m_k\geq2k+8$
for $k\geq2$. The constructed families give
the upper bound $\min\{4k+4,3k+8\}$; these inequalities do not determine
the exact minimum for $k\geq3$.
